
\textbf{Corollaire 4.4.} \uline{Sous les hypothèses de} 4.1, \uline{supposons} $S$ \uline{régulier}. \uline{Alors on a un isomorphisme} :
\[
K^{\bullet}(S)/J\xrightarrow{\sim}K^{\bullet}(P).
\]
Si $S$ est régulier, $G$ est aussi régulier, et $P$ également. On a donc $K^{\bullet}(S)\simeq K_\bullet(S)$, $K^{\bullet}(P)\simeq K_\bullet(P)$, d'où le résultat d'après 4.1.



% Exposé X local macros
\providecommand{\Filt}{\operatorname{Filt}}
\providecommand{\Coh}{\operatorname{Coh}}
\renewcommand{\Ob}{\operatorname{Ob}}
\providecommand{\longg}{\operatorname{long}}
\providecommand{\degtr}{\operatorname{deg.tr}}
\providecommand{\dimdot}{\operatorname{dim.}}
\providecommand{\Frac}{\operatorname{Frac}}
\providecommand{\fpqc}{\mathrm{fpqc}}
\providecommand{\Sch}{\operatorname{Sch}}
\providecommand{\Picf}{\operatorname{Pic}}
\providecommand{\Ens}{\operatorname{Ens}}
\providecommand{\im}{\operatorname{Im}}

